% Problem record op_42d3766b44e3a89b. This TeX file is derived from
% database/problems_json/op_42d3766b44e3a89b.json by scripts/sync-tex.mjs; edit the JSON record
% and rerun the script rather than editing this file.
\section{Quantum validity of the Zhang--Yeung inequality}
\label{sec:op_42d3766b44e3a89b}

\paragraph{Problem.}
Does the four-party Zhang--Yeung inequality hold for the von Neumann entropies of every finite-dimensional quantum state?

For a state $\rho_{ABCD}$, define

\begin{equation}
\begin{aligned}
S(R)&=-\operatorname{Tr}(\rho_R\log_2\rho_R),\\
I(A:B)&=S(A)+S(B)-S(AB),\\
I(A:B\mid C)&=S(AC)+S(BC)-S(C)-S(ABC).
\end{aligned}
\label{eq:42d3-1}
\end{equation}

Using the quantities in Eq.~\eqref{eq:42d3-1}, the proposed inequality is

\begin{equation}
2I(C:D)\leq I(A:B)+I(A:CD)+3I(C:D\mid A)+I(C:D\mid B)
\label{eq:42d3-2}
\end{equation}

Does Eq.~\eqref{eq:42d3-2} hold for every $\rho_{ABCD}$, without independence, separability, stabilizer, or holographic assumptions?

\subsection*{Status}

Unsolved

\subsection*{Source}

Bao, Cao, Walter, and Wang discuss this exact unrestricted four-party candidate in Section 4.2 \sourcecite{ref:42d3-bao15}{Bao15}. The statement is rewritten here to make its hypotheses and success criterion self-contained.

\subsection*{Progress}

\begin{itemize}
  \item For classical random variables, inequality~\eqref{eq:42d3-2} is the Zhang--Yeung non-Shannon inequality. It was proved by Zhang and Yeung and later treated systematically by Dougherty, Freiling, and Zeger \sourcecite{ref:42d3-zhang98}{Zhang98}\sourcecite{ref:42d3-dfz11}{DFZ11}.

  \item The 2025 Quantum Entropy Prover paper still identifies the search for additional unconstrained quantum entropy inequalities as open. Its Section V.A proves a related five-party inequality. Classically, an auxiliary-copy construction turns that precursor into a four-party non-Shannon inequality; the five-party quantum proof does not establish the required four-party quantum extension. \sourcecite{ref:42d3-huang25}{Huang25}

  \item A second important distinction is between failure to derive an inequality from strong subadditivity and an actual violating density operator. A linear-programming witness outside the basic entropy cone need not be realizable by a quantum state. Likewise, the known constrained quantum inequalities, which assume specified conditional mutual informations vanish, do not settle this unconstrained question. \sourcecite{ref:42d3-huang25}{Huang25}
\end{itemize}

\subsection*{References}

\begin{enumerate}[label={},leftmargin=4.8em,labelsep=0.5em]
  \footnotesize
  \item[\textup{[Zhang98]}]\label{ref:42d3-zhang98}
  Z. Zhang and R. W. Yeung, \emph{On Characterization of Entropy Function via Information Inequalities}, IEEE Transactions on Information Theory 44(4), 1440--1452 (1998). \href{https://doi.org/10.1109/18.681320}{doi:10.1109/18.681320}.

  \item[\textup{[DFZ11]}]\label{ref:42d3-dfz11}
  R. Dougherty, C. Freiling, and K. Zeger, \emph{Non-Shannon Information Inequalities in Four Random Variables}, arXiv:1104.3602 (2011). \href{https://arxiv.org/abs/1104.3602}{arXiv:1104.3602}.

  \item[\textup{[Bao15]}]\label{ref:42d3-bao15}
  N. Bao, C. Cao, M. Walter, and Z. Wang, \emph{Holographic entropy inequalities and gapped phases of matter}, arXiv:1507.05650v2 (2015), Section 4.2. This is the source for the unrestricted quantum candidate and the Ingleton distinction. \href{https://arxiv.org/html/1507.05650v2}{Full text}.

  \item[\textup{[Huang25]}]\label{ref:42d3-huang25}
  S.-L. Huang, T. Rippchen, and M. Berta, \emph{Quantum Entropy Prover}, arXiv:2501.16025v1, January 27, 2025, Sections V.A and V.C. \href{https://arxiv.org/html/2501.16025v1}{Full text}.
\end{enumerate}

\subsection*{Comment}

Retained as unresolved. Bao, Cao, Walter, and Wang provide an explicit source for the candidate, while the later entropy-cone work documents the broader gap in unconstrained quantum entropy inequalities.

Useful research target: Either prove the displayed inequality for unrestricted density operators or exhibit a finite-dimensional state with a rigorously certified negative difference between its right- and left-hand sides. Numerical minimization alone would be a search tool, not a proof of validity.

\subsection*{Field}

Quantum Communication

\subsection*{Topic}

Matrix and entropy inequalities

\subsection*{ID}

\texttt{op\_42d3766b44e3a89b}
